% Dokumentkopf in Anlehnung an die HSRM-Abschlussarbeitsvorlage

\documentclass[
  ngerman,
  12pt,
  a4paper,
  oneside,
  headings=normal
]{scrartcl}

\usepackage[ngerman]{babel}
\usepackage{fontspec}
\usepackage{microtype}
\usepackage{graphicx}
\usepackage{float}
\usepackage{xcolor}
\usepackage{booktabs}
\usepackage{longtable}
\usepackage{tabularx}
\usepackage{array}
\usepackage{enumitem}
\usepackage{ragged2e}
\usepackage{setspace}
\usepackage{csquotes}
\usepackage{amsmath}
\usepackage{amssymb}
\usepackage{tikz}
\usetikzlibrary{arrows.meta,positioning}
\usepackage[numbers,sort&compress]{natbib}
\usepackage{scrlayer-scrpage}

\definecolor{hsrmblue}{HTML}{005B82}
\definecolor{hsrmlight}{HTML}{EAF2F5}
\definecolor{hsrmgray}{HTML}{4B5563}
\definecolor{statusgreen}{HTML}{EAF6EF}
\definecolor{statusyellow}{HTML}{FFF7DC}

% Allgemeine Daten. Die eckig markierten Angaben sind vor Abgabe zu ersetzen.
\newcommand{\dokumenttyp}{Projekt-Expos\'e}
\newcommand{\titel}{Lokale Personen-Re-Identifikation als Kontextschicht f\"ur komplexe Innenraum- und IoT-Anwendungen}
\newcommand{\untertitel}{Konzeption, prototypische Umsetzung und Evaluation einer modularen Computer-Vision-Pipeline}
\newcommand{\autor}{[Name erg\"anzen]}
\newcommand{\matrikelnummer}{[Matrikelnummer erg\"anzen]}
\newcommand{\fachbereich}{Design Informatik Medien}
\newcommand{\studiengang}{[Studiengang erg\"anzen]}
\newcommand{\betreuung}{[Betreuung erg\"anzen]}
\newcommand{\abgabedatum}{25. August 2026}
\newcommand{\keywords}{Person Re-Identification, Computer Vision, Tracking, OSNet, YOLO, IoT, Privacy by Design}

\usepackage[
  colorlinks=true,
  linkcolor=hsrmblue,
  urlcolor=hsrmblue,
  citecolor=hsrmblue,
  bookmarksnumbered=true,
  pdftitle={\titel},
  pdfauthor={\autor},
  pdfkeywords={\keywords}
]{hyperref}

\KOMAoptions{parskip=half,headsepline=true}
\setlength{\parindent}{0pt}
\setlength{\textwidth}{150mm}
\setlength{\oddsidemargin}{10mm}
\setlength{\topmargin}{-4mm}
\setlength{\textheight}{235mm}
\setlength{\headheight}{18pt}
\setlength{\emergencystretch}{3em}

\clearpairofpagestyles
\ihead{\dokumenttyp}
\ohead{Person Re-Identification}
\cfoot{\pagemark}

\setcounter{tocdepth}{2}
\setcounter{secnumdepth}{2}

\newcommand{\code}[1]{\mbox{\texttt{#1}}}
\newcommand{\file}[1]{\mbox{\textsf{#1}}}
\newcommand{\statusdone}[1]{\colorbox{statusgreen}{\strut\small #1}}
\newcommand{\statusplan}[1]{\colorbox{statusyellow}{\strut\small #1}}

\newcolumntype{Y}{>{\RaggedRight\arraybackslash}X}
\newcolumntype{L}[1]{>{\RaggedRight\arraybackslash}p{#1}}

\tikzset{
  stage/.style={
    draw=hsrmblue,
    rounded corners=2pt,
    fill=hsrmlight,
    minimum width=30mm,
    minimum height=11mm,
    align=center,
    font=\small
  },
  flow/.style={-{Latex[length=2.2mm]}, thick, draw=hsrmblue}
}

\hypersetup{linktoc=all}
